\frfilename{mt112.tex}
\versiondate{20.2.05/20.8.08}

\def\chaptername{Ruang ukur}
\def\sectionname{Ruang ukur}
\copyrightdate{1994}

\newsection{112}

Sekarang, saya harap, kita telah siap untuk definisi besar yang kedua,
yakni definisi yang menjadi landasan seluruh pembahasan dalam karya ini.

\leader{112A}{Definisi} Sebuah {\bf ruang ukur} adalah tripel
$(X,\Sigma,\mu)$ dengan

\quad(i) $X$ suatu himpunan;

\quad(ii) $\Sigma$ suatu aljabar-$\sigma$ atas subhimpunan-subhimpunan $X$;

\quad(iii) $\mu:\Sigma\to[0,\infty]$ suatu fungsi sedemikian sehingga

\qquad($\alpha$) $\mu\emptyset=0$;

\qquad($\beta$) jika $\langle E_n\rangle_{n\in\Bbb N}$ adalah barisan
saling lepas di dalam $\Sigma$, maka
$\mu(\bigcup_{n\in\Bbb N}E_n)=\sum_{n=0}^{\infty}\mu E_n$.

\noindent Dalam konteks ini, anggota-anggota $\Sigma$ disebut himpunan
{\bf terukur}, dan $\mu$ disebut sebuah {\bf ukuran pada $X$}.

\leader{112B}{Catatan}\cmmnt{ {\bf (a) Penggunaan $\infty$} Dalam (iii)
pada definisi di atas, saya menetapkan bahwa $\mu$ adalah fungsi yang
mengambil nilai di `$[0,\infty]$', yaitu himpunan bilangan real
nonnegatif dengan ` $\infty$ ' ditambahkan. Saya menduga Anda telah
menjumpai berbagai penggunaan lambang $\infty$ dalam analisis; saya
harap Anda menyadari bahwa lambang ini mempunyai arti yang agak berbeda
dalam konteks yang berbeda, sehingga setiap kali digunakan perlu
ditetapkan kesepakatan yang jelas. ` $\infty$ pada ukuran' berkaitan
dengan gagasan panjang, luas, atau volume yang tak hingga. Operasi dasar
yang perlu kita lakukan terhadapnya ialah penjumlahan:
$\infty+a=a+\infty=\infty$ untuk setiap $a\in[0,\infty\roibr$ (yakni
setiap bilangan real $a\ge 0$), dan $\infty+\infty=\infty$. Dengan
demikian $[0,\infty]$ menjadi semigrup terhadap penjumlahan. Cukup aman
untuk menetapkan $\infty-a=\infty$ bagi setiap $a\in\Bbb R$; tetapi
kita sama sekali tidak boleh memberi makna pada rumus
$\infty-\infty$. Untuk perkalian, ternyata biasanya tepat untuk
menafsirkan $\infty\cdot\infty$, $a\cdot\infty$, dan $\infty\cdot a$
sebagai $\infty$ untuk $a>0$, sedangkan
$0\cdot\infty=\infty\cdot 0$ pada umumnya dapat diambil sebagai $0$.

Kita juga mempunyai urutan total alami pada $[0,\infty]$, dengan
menuliskan $a<\infty$ untuk setiap $a\in\coint{0,\infty}$. Ini memberi
pengertian supremum dan infimum bagi setiap subhimpunan (tak kosong)
dari $[0,\infty]$; dan sering kali tepat untuk menafsirkan
$\inf\emptyset$ sebagai $\infty$, tetapi saya akan berusaha menandai
kesepakatan khusus ini setiap kali relevan.
Kita juga mempunyai gagasan limit; jika
$\langle u_n\rangle_{n\in\Bbb N}$ adalah barisan di dalam
$[0,\infty]$, maka barisan itu konvergen ke $u\in[0,\infty]$ jika

\qquad {untuk setiap $v<u$ terdapat $n_0\in\Bbb N$ sedemikian sehingga
$v\le u_n$ untuk setiap $n\ge n_0$,}

\qquad {untuk setiap $v>u$ terdapat $n_0\in\Bbb N$ sedemikian sehingga
$v\ge u_n$ untuk setiap $n\ge n_0$.}

\noindent Tentu saja, jika $u=0$ atau $u=\infty$, salah satu syarat ini
akan terpenuhi secara hampa.

(Lihat juga \S135.)

\header{112Bb}{\bf (b)} Saya perlu menjelaskan dengan tegas apa yang
saya maksud dengan barisan `saling lepas': barisan
$\langle E_n\rangle_{n\in\Bbb N}$ disebut {\bf saling lepas} jika
tidak ada titik yang termasuk dalam lebih dari satu $E_n$, yakni jika
$E_m\cap E_n=\emptyset$ untuk setiap $m$, $n\in\Bbb N$ yang berbeda.
Perhatikan bahwa satu, bahkan banyak, di antara $E_n$ boleh merupakan
himpunan kosong.

Demikian pula, jika $\langle E_i\rangle_{i\in I}$ adalah keluarga
himpunan yang diindeks oleh sembarang himpunan $I$, keluarga itu disebut
{\bf saling lepas} jika $E_i\cap E_j=\emptyset$ untuk setiap $i$,
$j\in I$ yang berbeda.

\medskip

} {\bf (c)} Untuk menafsirkan syarat (iii-$\beta$) pada definisi di
atas, kita perlu menetapkan nilai jumlah $\sum_{n=0}^{\infty}u_n$ bagi
sembarang barisan $\langle u_n\rangle_{n\in\Bbb N}$ di dalam
$[0,\infty]$. \cmmnt{Cara alaminya ialah menetapkan
$\sum_{n=0}^{\infty}u_n=\lim_{n\to\infty}\sum_{m=0}^nu_m$,
dengan menggunakan definisi yang diuraikan dalam (a).} %end of comment
Jika salah satu $u_m$ sendiri tak hingga, \cmmnt{misalkan
$u_k=\infty$, maka $\sum_{m=0}^nu_m=\infty$ untuk setiap $n\ge k$,
sehingga tentu saja} $\sum_{n=0}^{\infty}u_n=\infty$. Jika semua $u_m$
berhingga, maka\cmmnt{, karena semuanya nonnegatif,} barisan jumlah
parsial $\langle\sum_{m=0}^nu_m\rangle_{n\in\Bbb N}$ bersifat monoton
tak-menurun, dan mempunyai limit berhingga
$\sum_{n=0}^{\infty}u_n\in\Bbb R$, atau divergen ke $\infty$; dalam
kasus terakhir ini kita kembali menafsirkan
$\sum_{n=0}^{\infty}u_n$ sebagai $\infty$.

\header{112Bd}{\bf (d)}\cmmnt{ Sekali lagi, contoh-contoh penting
ruang ukur harus menunggu sampai \S\S114 dan 115 di bawah. Namun, saya
dapat langsung menggambarkan satu kelas khusus ruang ukur yang
hendaknya selalu diingat, meskipun kelas ini tidak memberi gambaran
yang baik tentang bagian-bagian terpenting dan paling menarik dari
pokok bahasan ini.} Misalkan $X$ sembarang himpunan dan
$h:X\to[0,\infty]$ sembarang fungsi. Untuk setiap $E\subseteq X$,
tuliskan $\mu E=\sum_{x\in E}h(x)$. Untuk menafsirkan jumlah ini,
perhatikan bahwa tidak ada kesulitan jika $E$ berhingga (dengan
mengambil $\sum_{x\in\emptyset}h(x)=0$), sedangkan jika $E$ tak hingga
kita dapat mengambil
$\sum_{x\in E}h(x)=\sup\{\sum_{x\in I}h(x):I\subseteq E$ berhingga$\}$,
karena setiap $h(x)$ nonnegatif. \cmmnt{(Pada awalnya Anda mungkin
lebih suka memikirkan kasus $X=\Bbb N$, sehingga
$\sum_{n\in E}h(n)=\lim_{n\to\infty}\sum_{m\in E,m\le n}h(m)$; tetapi
saya harap sedikit pemikiran akan menunjukkan bahwa kasus umum, ketika
$X$ bahkan mungkin tak terhitung, sesungguhnya tidak lebih sulit.)}
Sekarang $(X,\Cal PX,\mu)$ adalah ruang ukur.

\cmmnt{Kita masih sangat jauh dari kesiapan untuk kosakata khusus yang
diperlukan guna menggambarkan berbagai jenis ruang ukur, tetapi ketika
saatnya tiba} saya akan menyebut ukuran semacam ini {\bf bertumpu pada
titik}.

Dua kasus khusus cukup sering muncul sehingga layak diberi nama. Jika
$h(x)=1$ untuk setiap $x$, maka $\mu E$ hanyalah banyaknya titik di
$E$ jika $E$ berhingga, dan bernilai $\infty$ jika $E$ tak hingga.
Saya akan menyebutnya {\bf ukuran cacah} pada $X$. Jika $x_0\in X$,
kita dapat menetapkan $h(x_0)=1$ dan $h(x)=0$ untuk
$x\in X\setminus\{x_0\}$; maka $\mu E$ bernilai $1$ jika $x_0\in E$,
dan $0$ untuk $E$ lainnya. Saya akan menyebutnya {\bf ukuran Dirac pada
$X$ yang terpusat di $x_0$}. Contoh sederhana lainnya ialah
$X=\Bbb N$, $h(n)=2^{-n-1}$ untuk setiap $n$; maka
$\mu X={1\over 2}+{1\over 4}+\ldots=1$.

\spheader 112Be Jika $(X,\Sigma,\mu)$ adalah ruang ukur\cmmnt{, maka
$\Sigma$ adalah domain fungsi $\mu$, dan $X$ adalah anggota terbesar
$\Sigma$. Karena itu, kita dapat memulihkan seluruh tripel
$(X,\Sigma,\mu)$ dari komponen terakhirnya, $\mu$. Tidak ada gunanya
melakukan ini pada tahap sekarang. Namun, kadang-kadang berguna untuk
memperkenalkan suatu ukuran tanpa segera memberi nama pada domainnya,
dan ketika melakukan ini} saya dapat mengatakan bahwa `$\mu$
{\bf mengukur} $E$' atau `$E$ {\bf diukur oleh} $\mu$' untuk menyatakan
bahwa $\mu E$ terdefinisi\cmmnt{, yakni bahwa $E$ termasuk dalam
aljabar-$\sigma$ $\dom\mu$}. \cmmnt{{\bf Peringatan!} Banyak penulis
menggunakan frasa `himpunan $\mu$-terukur' dengan arti yang agak berbeda
dari yang sedang saya bahas di sini.}

\leader{112C}{Sifat-sifat dasar ruang ukur} Misalkan
$(X,\Sigma,\mu)$ suatu ruang ukur.

(a) Jika $E$, $F\in\Sigma$ dan $E\cap F=\emptyset$, maka
$\mu(E\cup F)=\mu E+\mu F$.

(b) Jika $E$, $F\in\Sigma$ dan $E\subseteq F$, maka $\mu E\le \mu F$.

(c) $\mu(E\cup F)\le\mu E+\mu F$ untuk setiap $E$, $F\in\Sigma$.

(d) Jika $\langle E_n\rangle_{n\in\Bbb N}$ adalah sebarang barisan di
dalam $\Sigma$, maka
$\mu(\bigcup_{n\in\Bbb N}E_n)\le\sum_{n=0}^{\infty}\mu E_n$.

(e) Jika $\langle E_n\rangle_{n\in\Bbb N}$ adalah barisan
tak-menurun di dalam $\Sigma$ (yakni $E_n\subseteq E_{n+1}$ untuk
setiap $n\in\Bbb N$), maka

\Centerline{$\mu(\bigcup_{n\in\Bbb N}E_n)=\lim_{n\to\infty}\mu
E_n=\sup_{n\in\Bbb N}\mu E_n$.}

(f) Jika $\langle E_n\rangle_{n\in\Bbb N}$ adalah barisan tak-menaik
di dalam $\Sigma$ (yakni $E_{n+1}\subseteq E_{n}$ untuk setiap
$n\in\Bbb N$), dan jika salah satu $\mu E_n$ berhingga, maka

\Centerline{$\mu(\bigcap_{n\in\Bbb N}E_n)
=\lim_{n\to\infty}\mu E_n=\inf_{n\in\Bbb N}\mu E_n$.}

\proof{{\bf (a)} Tetapkan $E_0=E$, $E_1=F$, dan $E_n=\emptyset$
untuk $n\ge 2$; maka $\langle E_n\rangle_{n\in\Bbb N}$ adalah barisan
saling lepas di dalam $\Sigma$ dan
$\bigcup_{n\in\Bbb N}E_n=E\cup F$, sehingga

\Centerline{$\mu(E\cup
F)=\sum_{n=0}^{\infty}\mu E_n=\mu E+\mu F$}

\noindent (karena $\mu\emptyset=0$).

\medskip

{\bf (b)} $F\setminus E\in\Sigma$ (111Dc) dan
$\mu(F\setminus E)\ge 0$ (karena semua nilai $\mu$ termasuk dalam
$[0,\infty]$); jadi, dengan menggunakan (a),

\Centerline{$\mu F=\mu E+\mu(F\setminus E)\ge\mu E$.}

\medskip

{\bf (c)} Menurut (a),
$\mu(E\cup F)=\mu E+\mu(F\setminus E)$, dan menurut (b),
$\mu(F\setminus E)\le \mu F$.

\medskip

{\bf (d)} Tetapkan $F_0=E_0$ dan
$F_n=E_n\setminus\bigcup_{i<n}E_i$ untuk $n\ge 1$; maka
$\langle F_n\rangle_{n\in\Bbb N}$ adalah barisan saling lepas di dalam
$\Sigma$, $\bigcup_{n\in\Bbb N}F_n=\bigcup_{n\in\Bbb N}E_n$, dan
$F_n\subseteq E_n$ untuk setiap $n$. Menurut (b) tepat di atas,
$\mu F_n\le\mu E_n$ untuk setiap $n$; jadi

\Centerline{$\mu(\bigcup_{n\in\Bbb N}E_n)
=\mu(\bigcup_{n\in\Bbb N}F_n)=\sum_{n=0}^{\infty}\mu
F_n\le\sum_{n=0}^{\infty}\mu E_n$.}

\medskip

{\bf (e)} Tetapkan $F_0=E_0$ dan $F_n=E_n\setminus E_{n-1}$ untuk
$n\ge 1$; maka $\langle F_n\rangle_{n\in\Bbb N}$ adalah barisan
saling lepas di dalam $\Sigma$ dan
$\bigcup_{n\in\Bbb N}F_n=\bigcup_{n\in\Bbb N}E_n$. Akibatnya,
$\mu(\bigcup_{n\in\Bbb N}E_n)=\sum_{n=0}^{\infty}\mu F_n$. Namun,
induksi sederhana pada $n$, dengan menggunakan (a) untuk langkah
induksi, menunjukkan bahwa $\mu E_n=\sum_{m=0}^n\mu F_m$ untuk setiap
$n$. Jadi

\Centerline{$\sum_{n=0}^{\infty}\mu
F_n=\lim_{n\to\infty}\sum_{m=0}^n\mu F_m=\lim_{n\to\infty}\mu E_n$.}

\noindent Akhirnya,
$\lim_{n\to\infty}\mu E_n=\sup_{n\in\Bbb N}\mu E_n$ karena, menurut
(b), $\langle\mu E_n\rangle_{n\in\Bbb N}$ tak-menurun.

\medskip

{\bf (f)} Misalkan $\mu E_k<\infty$. Tetapkan
$F_n=E_k\setminus E_{k+n}$ untuk $n\in\Bbb N$ dan
$F=\bigcup_{n\in\Bbb N}F_n$; maka
$\langle F_n\rangle_{n\in\Bbb N}$ adalah barisan tak-menurun di dalam
$\Sigma$, sehingga menurut (e) tepat di atas,
$\mu F=\lim_{n\to\infty}\mu F_n$. Selain itu,
$\mu F_n+\mu E_{k+n}=\mu E_k$; karena $\mu E_k<\infty$, kita dapat
dengan aman menuliskan $\mu F_n=\mu E_k-\mu E_{k+n}$, sehingga

\Centerline{$\mu F=\lim_{n\to\infty}(\mu E_k-\mu
E_{k+n})=\mu E_k-\lim_{n\to\infty}\mu E_n$.}

\noindent Selanjutnya, $F\subseteq E_k$, sehingga
$\mu F+\mu(E_k\setminus F)=\mu E_k$, dan (sekali lagi karena
$\mu E_k$ berhingga) $\mu F=\mu E_k-\mu(E_k\setminus F)$. Jadi haruslah
$\mu(E_k\setminus F)=\lim_{n\to\infty}\mu E_n$. Namun,
$E_k\setminus F$ tepat sama dengan $\bigcap_{n\in\Bbb N}E_n$.

Akhirnya, $\lim_{n\to\infty}\mu E_n=\inf_{n\in\Bbb N}\mu E_n$
karena $\langle\mu E_n\rangle_{n\in\Bbb N}$ tak-menaik.
}%end of proof of 112C

\medskip

\cmmnt{\noindent{\bf Catatan} Perhatikan bahwa dalam (f) di atas,
syarat $\inf_{n\in\Bbb N}\mu E_n<\infty$ sangatlah penting.
Konstruksi dalam 112Bd sudah cukup untuk menunjukkan hal ini. Ambil
$X=\Bbb N$ dan misalkan $\mu$ ukuran cacah pada $X$. Tetapkan
$E_n=\{i:i\in\Bbb N,\,i\ge n\}$ untuk setiap $n$. Maka
$E_{n+1}\subseteq E_n$ untuk setiap $n$, tetapi

\Centerline{$\mu(\bigcap_{n\in\Bbb N}E_n)
=\mu\emptyset=0<\infty=\lim_{n\to\infty}\mu E_n$.}
}%end of comment


\leader{112D}{Himpunan terabaikan} Misalkan $(X,\Sigma,\mu)$ sembarang
ruang ukur.

\header{112Da}{\bf (a)} Suatu himpunan $A\subseteq X$ disebut
{\bf terabaikan} (atau {\bf nol}) jika terdapat $E\in\Sigma$
sedemikian sehingga $A\subseteq E$ dan $\mu E=0$.
(Jika mungkin timbul keraguan tentang ukuran mana yang digunakan, saya
akan menuliskan {\bf terabaikan-$\mu$}.)

\header{112Db}{\bf (b)} Misalkan $\Cal N$ keluarga subhimpunan
terabaikan dari $X$. Maka (i) $\emptyset\in\Cal N$ (ii) jika
$A\subseteq B\in\Cal N$, maka $A\in\Cal N$ (iii) jika
$\langle A_n\rangle_{n\in\Bbb N}$ adalah sembarang barisan di dalam
$\Cal N$, maka $\bigcup_{n\in\Bbb N}A_n\in\Cal N$.
\prooflet{\Prf\ (i) $\mu(\emptyset)=0$. (ii) Terdapat $E\in\Sigma$
sedemikian sehingga $\mu E=0$ dan $B\subseteq E$; sekarang
$A\subseteq E$. (iii) Untuk setiap $n\in\Bbb N$, pilih
$E_n\in\Sigma$ sedemikian sehingga $A_n\subseteq E_n$ dan $\mu E_n=0$.
Sekarang $E=\bigcup_{n\in\Bbb N}E_n\in\Sigma$ dan
$\bigcup_{n\in\Bbb N}A_n\subseteq\bigcup_{n\in\Bbb N}E_n$, serta
$\mu(\bigcup_{n\in\Bbb N}E_n)\le\sum_{n=0}^{\infty}\mu E_n$ menurut
112Cd, sehingga $\mu(\bigcup_{n\in\Bbb N}E_n)=0$. \Qed}

Saya akan menyebut $\Cal N$ sebagai {\bf ideal nol} dari ukuran $\mu$.
(Keluarga himpunan yang memenuhi syarat-syarat (i)-(iii) di sini
disebut {\bf ideal-$\sigma$} himpunan.)

\header{112Dc}{\bf (c)} Suatu himpunan $A\subseteq X$ disebut
{\bf koterabaikan} jika $X\setminus A$ terabaikan\cmmnt{; yakni,
terdapat himpunan terukur $E\subseteq A$ sedemikian sehingga
$\mu(X\setminus E)=0$}. Perhatikan bahwa (i) $X$ koterabaikan (ii) jika
$A\subseteq B\subseteq X$ dan $A$ koterabaikan, maka $B$
koterabaikan (iii) jika $\langle A_n\rangle_{n\in\Bbb N}$ adalah
barisan himpunan koterabaikan, maka $\bigcap_{n\in\Bbb N}A_n$
koterabaikan.

\header{112Dd}{\bf (d)}\cmmnt{ Berguna, dan lazim, untuk menggunakan
bahasa yang agak informal mengenai himpunan terabaikan.} Jika $P(x)$
adalah suatu pernyataan yang berlaku bagi anggota $x$ dari himpunan
$X$, kita mengatakan bahwa

\Centerline{`$P(x)$ untuk hampir setiap $x\in X$'}

\noindent atau

\Centerline{`$P(x)$ a.e.\ $(x)$'}

\noindent atau

\Centerline{`$P$ hampir di mana-mana',\quad `$P$ a.e.'}

\noindent atau\cmmnt{, jika ukuran yang digunakan tampaknya perlu
dinyatakan,}

\Centerline{`$P(x)$ untuk $\mu$-hampir setiap $x$',
\quad`$P(x)\,\mu$-a.e.$(x)$',
\quad`$P\,\mu$-a.e.',}

\noindent untuk menyatakan bahwa

\Centerline{$\{x:x\in X,\,P(x)\}$}

\noindent koterabaikan di dalam $X$\cmmnt{, yakni bahwa

\Centerline{$\{x:x\in X,\,P(x)$ salah$\}$}

\noindent terabaikan}. Jadi\cmmnt{, misalnya,} jika
$f:X\to\Bbb R$ suatu fungsi, `$f>0$ a.e.' berarti bahwa
$\{x:f(x)\le 0\}$ terabaikan.

\cmmnt{\spheader 112De Frasa `{\bf hampir pasti}' ({\bf a.s.}) dan
`{\bf presque partout}' ({\bf p.p.}) juga digunakan untuk `hampir di
mana-mana'.
}%end of comment

\cmmnt{\header{112Df}{\bf (f)} Saya perlu menarik perhatian Anda pada
kenyataan bahwa, menurut definisi saya, suatu himpunan terabaikan tidak
harus terukur, walaupun harus termuat dalam suatu himpunan terukur yang
terabaikan. (Ruang ukur yang setiap himpunan terabaikannya terukur
disebut {\bf lengkap}. Saya akan kembali ke persoalan ini dalam \S211.)
}%end of comment

\spheader 112Dg Jika $f$ dan $g$ adalah fungsi bernilai real yang
didefinisikan pada subhimpunan-subhimpunan koterabaikan dari suatu
ruang ukur, saya akan menuliskan $f\eae g$, $f\leae g$, atau
$f\geae g$ untuk masing-masing menyatakan

\Centerline{$f=g$ a.e., yakni,
$\{x:x\in\dom(f)\cap\dom(g)$, $f(x)=g(x)\}$ koterabaikan,}

\Centerline{$f\le g$ a.e., yakni,
$\{x:x\in\dom(f)\cap\dom(g)$, $f(x)\le g(x)\}$ koterabaikan,}

\Centerline{$f\ge g$ a.e., yakni,
$\{x:x\in\dom(f)\cap\dom(g)$, $f(x)\ge g(x)\}$ koterabaikan.}

\exercises{
\leader{112X}{Latihan dasar $\pmb{>}$(a)}
%\sqheader 112Xa
Misalkan $(X,\Sigma,\mu)$ suatu ruang ukur. Tunjukkan bahwa (i)
$\mu(E\cup F)+\mu(E\cap F)=\mu E+\mu F$ (ii)
$\mu(E\cup F\cup G)+\mu(E\cap F)+\mu(E\cap G)+\mu(F\cap G)=\mu E+\mu
F+\mu G+\mu(E\cap F\cap G)$ untuk semua $E$, $F$, $G\in\Sigma$.
Perumum hasil-hasil ini untuk barisan himpunan yang lebih panjang.
(Anda mungkin lebih suka memulai dengan kasus ketika $\mu E$, $\mu F$,
dan $\mu G$ semuanya berhingga. Namun, saya harap Anda dapat menemukan
argumen yang menangani kasus umum.)
%112C

\sqheader 112Xb Misalkan $(X,\Sigma,\mu)$ suatu ruang ukur dan
$\langle E_n\rangle_{n\in\Bbb N}$ suatu barisan di dalam $\Sigma$.
Tunjukkan bahwa

\Centerline{$\mu(\bigcup_{n\in\Bbb N}\bigcap_{m\ge n}E_m)
\le\liminf_{n\to\infty}\mu E_n$.}
%112C

\spheader 112Xc Misalkan $(X,\Sigma,\mu)$ suatu ruang ukur dan
$E$, $F\in\Sigma$; andaikan $\mu E<\infty$. Tunjukkan bahwa
$\mu(F\symmdiff E)=\mu F-\mu E+2\mu(E\setminus F)$.
%112C

\spheader 112Xd Misalkan $(X,\Sigma,\mu)$ ruang ukur dan
$\sequencen{E_n}$ barisan himpunan terukur dengan
$\mu(\bigcup_{n\in\Bbb N}E_n)<\infty$. (i) Tunjukkan bahwa
$\limsup_{n\to\infty}\mu E_n
\le\mu(\bigcap_{n\in\Bbb N}\bigcup_{m\ge n}E_m)$. (ii) Tunjukkan
bahwa jika $\bigcap_{n\in\Bbb N}\bigcup_{m\ge n}E_m=E
=\bigcup_{n\in\Bbb N}\bigcap_{m\ge n}E_m$, maka
$\lim_{n\to\infty}\mu E_n$ ada dan sama dengan $\mu E$.
%112Xb 112C

\sqheader 112Xe Misalkan $(X,\Sigma,\mu)$ suatu ruang ukur, dan
$\eusm F$ himpunan fungsi bernilai real yang domainnya merupakan
subhimpunan koterabaikan dari $X$. (i) Tunjukkan bahwa
$\{(f,g):f,\,g\in\eusm F,\,f\leae g\}$ dan
$\{(f,g):f,\,g\in\eusm F,\,f\geae g\}$ adalah relasi refleksif dan
transitif pada $\eusm F$, dan masing-masing merupakan invers dari yang
lain. (ii) Tunjukkan bahwa
$\{(f,g):f,\,g\in\eusm F,\,f\eae g\}$ adalah irisan keduanya dan
merupakan relasi ekuivalensi pada $\eusm F$.
%112D

\spheader 112Xf Misalkan $(X,\Sigma,\mu)$ suatu ruang ukur, $Y$ suatu
himpunan, dan $\phi:X\to Y$ suatu fungsi. Tetapkan
$\Tau=\{F:F\subseteq Y$, $\phi^{-1}[F]\in\Sigma\}$ dan
$\nu F=\mu\phi^{-1}[F]$ untuk $F\in\Tau$. Tunjukkan bahwa $\nu$ adalah
ukuran pada $Y$. ($\nu$ disebut {\bf ukuran citra} pada $Y$, dan
umumnya akan saya lambangkan dengan $\mu\phi^{-1}$.)
%111Xc

\leader{112Y}{Latihan lanjutan (a)}
%\spheader 112Ya
Misalkan $X$ suatu himpunan dan $\Sigma$ suatu aljabar-$\sigma$ atas
subhimpunan-subhimpunan $X$. Misalkan $\mu_1$ dan $\mu_2$ dua ukuran
pada $X$, keduanya dengan domain $\Sigma$. Tetapkan

\Centerline{$\mu E
=\inf\{\mu_1(E\cap F)+\mu_2(E\setminus F):F\in\Sigma\}$}

\noindent untuk setiap $E\in\Sigma$. Tunjukkan bahwa $\mu$ adalah
ukuran pada $X$, dan merupakan ukuran terbesar, dengan domain $\Sigma$,
sedemikian sehingga $\mu E\le\min(\mu_1E,\mu_2E)$ untuk setiap
$E\in\Sigma$.
%problem with double use of symbol  \mu , here and in next exercise,
%see correspondence with Wallace Thompson

\spheader 112Yb Misalkan $X$ suatu himpunan dan $\Sigma$ suatu
aljabar-$\sigma$ atas subhimpunan-subhimpunan $X$. Misalkan $\mu_1$
dan $\mu_2$ dua ukuran pada $X$, keduanya dengan domain $\Sigma$.
Tetapkan

\Centerline{$\mu E
=\sup\{\mu_1(E\cap F)+\mu_2(E\setminus F):F\in\Sigma\}$}

\noindent untuk setiap $E\in\Sigma$. Tunjukkan bahwa $\mu$ adalah
ukuran pada $X$, dan merupakan ukuran terkecil, dengan domain $\Sigma$,
sedemikian sehingga $\mu E\ge\max(\mu_1E,\mu_2E)$ untuk setiap
$E\in\Sigma$.

\spheader 112Yc Misalkan $X$ suatu himpunan dan $\Sigma$ suatu
aljabar-$\sigma$ atas subhimpunan-subhimpunan $X$.

\quad(i) Andaikan $\Nu=\{\nu_0,\ldots,\nu_n\}$ adalah keluarga
ukuran-ukuran pada $X$,
semuanya dengan domain $\Sigma$. Tetapkan

\Centerline{$\mu E=\inf\{\sum_{i=0}^n\nu_iF_i:F_0,\ldots,F_n\in\Sigma$,
$E\subseteq\bigcup_{i\le n}F_i\}$}

\noindent untuk $E\in\Sigma$. Tunjukkan bahwa $\mu$ adalah ukuran
pada $X$.

\quad(ii) Misalkan $\Nu$ suatu keluarga tak kosong yang terdiri atas
ukuran-ukuran pada $X$, semuanya dengan domain $\Sigma$. Tetapkan

$$\eqalign{\mu E=\inf\{\sum_{n=0}^{\infty}\nu_nF_n:
&\sequencen{\nu_n}\text{ adalah barisan di dalam }\Nu,\cr
&\sequencen{F_n}\text{ adalah barisan di dalam }\Sigma,
\,E\subseteq\bigcup_{n\in\Bbb N}F_n\}\cr}$$

\noindent untuk $E\in\Sigma$. Tunjukkan bahwa $\mu$ adalah ukuran
pada $X$.

\quad(iii) Misalkan $\Nu$ suatu keluarga tak kosong yang terdiri atas
ukuran-ukuran pada $X$, semuanya dengan domain $\Sigma$, dan andaikan terdapat
$\nuprime\in\Nu$ sedemikian sehingga $\nuprime X<\infty$. Tetapkan

\Centerline{$\mu E=\inf\{\sum_{i=0}^n\nu_iF_i:n\in\Bbb N,
  \,\nu_0,\ldots,\nu_n\in\Nu,\,F_0,\ldots,F_n\in\Sigma,
  E\subseteq\bigcup_{i\le n}F_i\}$}

\noindent untuk $E\in\Sigma$. Tunjukkan bahwa $\mu$ adalah ukuran
pada $X$.

\quad(iv) Andaikan, dalam (iii), bahwa $\Nu$ terarah ke bawah, yakni
untuk setiap $\nu_1$, $\nu_2\in\Nu$ terdapat $\nu\in\Nu$ sedemikian
sehingga $\nu E\le\min(\nu_1E,\nu_2E)$ untuk setiap $E\in\Sigma$.
Tunjukkan bahwa $\mu E=\inf_{\nu\in\Nu}\nu E$ untuk setiap
$E\in\Sigma$.

\quad(v) Tunjukkan bahwa dalam semua kasus (i)-(iii), ukuran yang
dikonstruksi adalah ukuran terbesar $\mu$ dengan domain $\Sigma$
sedemikian sehingga $\mu E\le\inf_{\nu\in\Nu}\nu E$ untuk setiap
$E\in\Sigma$.

\spheader 112Yd Misalkan $X$ suatu himpunan dan $\Sigma$ suatu
aljabar-$\sigma$ atas subhimpunan-subhimpunan $X$. Misalkan $\Nu$
suatu keluarga tak kosong yang terdiri atas ukuran-ukuran pada $X$,
semuanya dengan domain $\Sigma$. Tetapkan

$$\eqalign{\mu E=\sup\{\sum_{i=0}^n\nu_iF_i:n\in\Bbb N,
  &\,\nu_0,\ldots,\nu_n\in\Nu,\cr
  &F_0,\ldots,F_n\text{ adalah subhimpunan saling lepas dari }E
    \text{ yang termasuk dalam }\Sigma\}\cr}$$

\noindent untuk $E\in\Sigma$. (i) Tunjukkan bahwa

$$\eqalign{\mu E=\sup\{\sum_{n=0}^{\infty}\nu_nF_n:
&\sequencen{\nu_n}\text{ adalah barisan di dalam }\Nu,\cr
&\sequencen{F_n}\text{ adalah barisan saling lepas di dalam }\Sigma,
\,\bigcup_{n\in\Bbb N}F_n\subseteq E\}\cr}$$

\noindent untuk setiap $E\in\Sigma$.
(ii) Tunjukkan bahwa $\mu$ adalah ukuran pada $X$, dan merupakan ukuran
terkecil, dengan domain $\Sigma$, sedemikian sehingga
$\mu E\ge\sup_{\nu\in\Nu}\nu E$ untuk setiap $E\in\Sigma$. (iii)
Sekarang andaikan $\Nu$ terarah ke atas, yakni untuk setiap $\nu_1$,
$\nu_2\in\Nu$ terdapat $\nu\in\Nu$ sedemikian sehingga
$\nu E\ge\max(\nu_1E,\nu_2E)$ untuk setiap $E\in\Sigma$. Tunjukkan
bahwa $\mu E=\sup_{\nu\in\Nu}\nu E$ untuk setiap $E\in\Sigma$.

\spheader 112Ye Misalkan $(X,\Sigma,\mu)$ suatu ruang ukur dan
$\sequencen{E_n}$ suatu barisan himpunan terukur. Untuk setiap
$k\in\Bbb N$, tetapkan
$H_k=\{x:x\in X,\,\#(\{n:x\in E_n\})\ge k\}$, yaitu himpunan titik
yang termasuk dalam $E_n$ untuk $k$ atau lebih nilai $n$. (i) Tunjukkan
bahwa setiap $H_k$ terukur. (ii) Tunjukkan bahwa
$\sum_{k=1}^{\infty}\mu H_k=\sum_{n=0}^{\infty}\mu E_n$.
\Hint{mulailah dengan kasus ketika $E_n=\emptyset$ untuk $n\ge n_0$.}
(iii) Tunjukkan bahwa jika $\sum_{n=0}^{\infty}\mu E_n$ berhingga,
maka hampir setiap titik di $X$ termasuk hanya dalam berhingga banyak
$E_n$, dan $\sum_{n=0}^{\infty}\mu E_n=\sum_{k=0}^{\infty}k\mu G_k$,
dengan

\Centerline{$G_k=H_k\setminus H_{k+1}=\{x:\#(\{n:x\in E_n\})=k\}$.}

\spheader 112Yf Misalkan $X$ suatu himpunan dan $\mu$, $\nu$ dua
ukuran pada $X$ dengan domain masing-masing $\Sigma$ dan $\Tau$.
Tetapkan $\Lambda=\Sigma\cap\Tau$ dan definisikan
$\lambda:\Lambda\to[0,\infty]$ dengan menetapkan
$\lambda E=\mu E+\nu E$ untuk setiap $E\in\Lambda$. Tunjukkan bahwa
$(X,\Lambda,\lambda)$ adalah ruang ukur.
%+
}%end of exercises

\endnotes{
\Notesheader{112} Perhitungan dalam hasil-hasil seperti 112Ca-112Cc,
112Xa, dan 112Xc, yang hanya melibatkan berhingga banyak himpunan,
berlaku bagi setiap konsep ukuran yang aditif; Anda mungkin pernah
menjumpainya dalam teori probabilitas dasar, tetapi tentu saja sekarang
saya meminta Anda mempertimbangkan pula kemungkinan bahwa satu atau
lebih himpunan mempunyai ukuran $\infty$. Saya harap Anda akan mendapati
bahwa hasil-hasil ini sepenuhnya alamiah dan tidak mengejutkan. Dalam
konteks ini saya menyarankan diagram Venn; hasil semacam ini yang hanya
melibatkan berhingga banyak himpunan terukur dan hanya {\it penjumlahan},
tanpa {\it pengurangan}, berlaku secara umum jika dan hanya jika hasil
itu berlaku bagi luas bangun geometris sederhana pada bidang. Syarat
`$\mu E<\infty$' dalam 112Xc diperlukan hanya karena kita mengurangkan
$\mu E$; hasil aditif yang bersesuaian,
$\mu(F\symmdiff E)+\mu E=\mu F+2\mu(E\setminus F)$, berlaku untuk semua
$E$ dan $F$ yang terukur.
Tentu saja, ketika barisan himpunan mulai terlibat, kita perlu sedikit
lebih berhati-hati; hasil-hasil 112Cd-112Cf adalah hasil dasar yang
perlu dipelajari. Namun, menurut saya satu-satunya jebakan terdapat pada
syarat `salah satu $\mu E_n$ berhingga' dalam 112Cf. Sebagaimana dicatat
dalam catatan pada akhir 112C, syarat ini sangat penting, dan untuk
barisan menurun yang terdiri atas himpunan-himpunan terukur, ukuran
himpunan limit mungkin benar-benar lebih
kecil daripada limit ukuran-ukurannya, walaupun hal ini hanya dapat
terjadi ketika limit terakhir tak hingga.
}%end of notes

\discrpage
